% =====================================================================
%  CHAPTER 1: 绪论
% =====================================================================
\chapter{绪论}
\setcounter{chapter}{1}
\chapterintro{
掌握社会医学的概念、研究对象和研究内容；了解社会医学的任务与发展历程；理解社会医学的特色理论与创新观点。
}

% ----- 1.1 社会医学概述 -----
\section{社会医学概述}

\begin{defbox}
\textbf{社会医学（Social Medicine）}是研究社会因素与健康及疾病之间相互联系及其规律的一门科学，是医学与社会学、经济学、管理学等社会科学相互渗透融合形成的一门交叉学科。它从社会宏观视角出发，研究社会因素对人群健康的影响，探讨保障人群健康的社会策略和方法。
\end{defbox}

\subsection{社会医学的研究对象}

社会医学以\keyword{社会人群的健康}为研究对象，主要关注：
\begin{enumerate}[leftmargin=*]
    \item \imp{社会因素与人群健康的相互作用关系}——研究社会政治、经济、文化、行为、环境等因素对人群健康的影响
    \item \imp{社会卫生状况及人群健康状况}——评价和描述人群健康状况及其变动规律
    \item \imp{改善社会卫生状况的策略与措施}——提出保护人群健康、提高健康水平的社会卫生对策
    \item \imp{社会医学研究方法}——运用社会学、流行病学、统计学等多学科方法开展研究
\end{enumerate}

\subsection{社会医学的研究内容}

\begin{tcolorbox}[colback=light, colframe=main, boxrule=0.5pt, arc=4pt, title=\textbf{社会医学的三大研究内容}, breakable]
\begin{enumerate}[leftmargin=*, itemsep=3pt]
    \item \imp{研究社会卫生状况}——主要是人群健康状况。通过社会卫生调查，评价人群健康水平及其变动规律，发现主要社会卫生问题及高危人群。核心指标包括：死亡率、发病率、期望寿命、伤残调整寿命年（DALY）等。
    \item \imp{研究影响人群健康的因素}——特别是社会制度、经济状况、文化因素、人口发展、生活劳动条件、行为生活方式及卫生服务等社会因素对健康的影响。揭示社会因素与健康之间的因果关系及其强度。
    \item \imp{研究社会卫生策略与措施}——针对已发现的社会卫生问题及其原因，制定并实施有效的社会卫生策略与措施，包括卫生政策、卫生规划、卫生服务体系改革、健康促进等。
\end{enumerate}
\end{tcolorbox}

\begin{keybox}
社会医学与\textbf{预防医学}、\textbf{社区医学}、\textbf{卫生管理学}、\textbf{医学社会学}等学科既有联系又有区别：
\begin{itemize}[leftmargin=*]
    \item \textbf{预防医学}侧重于生物、物理、化学因素对健康的影响及其预防
    \item \textbf{社会医学}侧重于社会因素对健康的影响及社会性防治策略
    \item \textbf{社区医学}侧重于社区范围内的卫生服务和健康管理
    \item \textbf{卫生管理学}侧重于卫生系统的组织、管理和政策实施
    \item \textbf{医学社会学}从社会学角度研究医学现象、医疗组织和医患关系
\end{itemize}
\end{keybox}

% ----- 1.2 社会医学的任务 -----
\section{社会医学的任务}

\begin{enumerate}[leftmargin=*]
    \item \imp{倡导积极的健康观}——推动"生物-心理-社会医学模式"在医疗卫生实践中的贯彻，促进健康观念的更新
    \item \imp{发现和评价社会卫生问题}——通过社会卫生调查和研究，发现主要的社会卫生问题，评价其严重程度和影响范围
    \item \imp{揭示社会因素与健康的因果关系}——阐明社会因素对人群健康的正面和负面影响及其作用机制
    \item \imp{制定社会卫生策略和措施}——针对社会卫生问题提出科学、可行的干预对策，改善社会卫生状况
    \item \imp{培养社会医学专业人才}——提高医疗卫生工作者的社会医学素养，适应新医学模式要求
\end{enumerate}

% ----- 1.3 社会医学的发展 -----
\section{社会医学的发展历程}

\subsection{国外社会医学发展}

\begin{table}[htbp]
\centering
\caption{社会医学发展史上的重要人物与贡献}
\begin{tabularx}{\textwidth}{lX}
\hline
\textbf{人物/时期} & \textbf{主要贡献} \\
\hline
\textbf{古希腊时期} & 希波克拉底（Hippocrates）提出环境与健康的关系，关注空气、水、土壤和居住条件对健康的影响 \\
\textbf{18-19世纪} & 工业革命时期，城市化带来严重公共卫生问题，推动卫生改革运动 \\
\textbf{法国} Rochoux（1838）& 首次提出"社会医学"概念 \\
\textbf{法国} Guérin（1848）& 将社会医学分为社会生理学、社会病理学、社会卫生学、社会治疗学四部分 \\
\textbf{德国} Neumann（1813-1908）& 强调"医学是一门社会科学"，健康是社会问题 \\
\textbf{德国} Virchow（1821-1902）& 提出"医学科学的核心是社会科学"、"政治是更大范围的医学"，开展斑疹伤寒调查 \\
\textbf{德国} Grotjahn（1869-1931）& 现代实验社会医学奠基人，1920年创办柏林大学社会卫生学系 \\
\textbf{英国} 19世纪 & Chadwick推动公共卫生法案；Simon开展社会卫生调查 \\
\textbf{WHO（1948起）}& 提出健康定义；推动"人人享有卫生保健"全球战略（1978，Alma-Ata宣言） \\
\hline
\end{tabularx}
\end{table}

\begin{keybox}
\textbf{社会医学发展的三个时期：}
\begin{enumerate}[leftmargin=*]
    \item \textbf{萌芽期（19世纪以前）}：以临床医学和环境卫生学为基础，关注环境与疾病的关系
    \item \textbf{创立期（19世纪）}：社会医学概念正式提出，卫生改革运动兴起，关注社会因素对人群健康的影响
    \item \textbf{发展期（20世纪至今）}：社会医学成为独立学科，进入大学教育体系，WHO推动全球卫生战略
\end{enumerate}
\end{keybox}

\subsection{我国社会医学发展}

\begin{enumerate}[leftmargin=*]
    \item \textbf{古代}：中医经典《黄帝内经》即重视环境、情志与社会因素对人体健康的影响，体现"天人合一"的整体医学思想
    \item \textbf{近代（20世纪初）}：1910年设立卫生检疫机构；1922年《社会医学》概念传入中国；1941年陈志潜等在河北定县开展农村卫生实验区，创建"定县模式"，被誉为"中国社区医学之父"
    \item \textbf{新中国成立后（1949-）}：1952年设立卫生系，开展社会卫生教学；引进苏联保健组织学；开展爱国卫生运动
    \item \textbf{改革开放后（1980-）}：1980年卫生部发布《关于加强社会医学与卫生管理学教学科研工作的通知》；1984年第一本《社会医学》教材出版；各医学院校设立社会医学教研室
    \item \textbf{21世纪以来}：社会医学与卫生事业管理学科蓬勃发展；健康中国战略实施；社区卫生服务和全科医学快速发展
\end{enumerate}

% ----- 1.4 社会医学的特色理论与创新观点 -----
\section{社会医学的特色理论与创新观点}

\begin{tcolorbox}[colback=light, colframe=main, boxrule=0.5pt, arc=4pt, title=\textbf{社会医学的特色理论}, breakable]
\begin{enumerate}[leftmargin=*]
    \item \imp{社会健康观}——健康不仅是生物学现象，更是社会现象。人群健康水平受到政治、经济、文化、环境等社会因素的综合影响
    \item \imp{社会因素决定论}——社会因素是影响人群健康的根本原因，生物因素和行为因素受社会因素制约
    \item \imp{健康公平理论}——关注不同社会阶层、地区、人群之间健康水平的差异，强调社会公正和健康公平
    \item \imp{整体健康观}——健康是社会、心理、生物因素相互作用的综合体现，需从多维度认识健康
    \item \imp{大卫生观}——健康不仅是卫生部门的责任，需要全社会参与，强调跨部门协作
    \item \imp{健康社会决定因素理论}——WHO倡导的健康社会决定因素框架，强调"从根源上解决健康不公平"
\end{enumerate}
\end{tcolorbox}

\begin{keybox}
\textbf{社会医学的创新观点：}
\begin{itemize}[leftmargin=*]
    \item \keyword{将健康融入所有政策（HiAP）}：健康不仅是卫生部门的职责，所有公共政策都应考虑对健康的影响
    \item \keyword{健康共治}：政府、社会、个人共同参与健康治理
    \item \keyword{全生命周期健康管理}：从出生到死亡的全过程健康维护
    \item \keyword{健康中国战略}：将人民健康放在优先发展的战略地位
\end{itemize}
\end{keybox}

% ----- 复习题 -----
\section{复习题}

\subsection*{一、选择题}
\begin{enumerate}[leftmargin=*, itemsep=3pt]
    \item 社会医学是研究什么的学科？
    \begin{enumerate}
        \item[(A)] 社会因素与健康及疾病之间相互联系及其规律
        \item[(B)] 生物因素与疾病的关系
        \item[(C)] 环境卫生与健康的关系
        \item[(D)] 医疗技术与疾病治疗的关系
    \end{enumerate}
    \item 首次提出"社会医学"概念的是哪位学者？
    \begin{enumerate}
        \item[(A)] Virchow
        \item[(B)] Neumann
        \item[(C)] Rochoux
        \item[(D)] Grotjahn
    \end{enumerate}
    \item 社会医学的三大研究内容不包括：
    \begin{enumerate}
        \item[(A)] 研究社会卫生状况
        \item[(B)] 研究影响人群健康的因素
        \item[(C)] 研究临床治疗方法
        \item[(D)] 研究社会卫生策略与措施
    \end{enumerate}
    \item 被誉为"现代实验社会医学奠基人"的是谁？
    \begin{enumerate}
        \item[(A)] Virchow
        \item[(B)] Neumann
        \item[(C)] Guérin
        \item[(D)] Grotjahn
    \end{enumerate}
\end{enumerate}

\subsection*{二、名词解释}
\begin{enumerate}[leftmargin=*, itemsep=3pt]
    \item \textbf{社会医学（Social Medicine）}
    \item \textbf{大卫生观}
    \item \textbf{将健康融入所有政策（HiAP）}
\end{enumerate}

\subsection*{三、简答题}
\begin{enumerate}[leftmargin=*, itemsep=3pt]
    \item 简述社会医学的研究对象与研究内容。
    \item 试述社会医学的基本任务。
    \item 简述社会医学的特色理论观点。
    \item 简述我国社会医学发展的重要阶段及标志性事件。
\end{enumerate}

\subsection*{参考答案要点}

\subsubsection*{一、选择题}
1. (A)；2. (C)；3. (C)；4. (D)

\subsubsection*{二、名词解释}
\begin{enumerate}[leftmargin=*, itemsep=3pt]
    \item \textbf{社会医学}：研究社会因素与健康及疾病之间相互联系及其规律的一门交叉学科，从社会宏观视角研究社会因素对人群健康的影响，探讨保障人群健康的社会策略和方法。
    \item \textbf{大卫生观}：健康不仅是卫生部门的责任，需要全社会参与和跨部门协作，将健康纳入所有政策领域的宏观卫生理念。
    \item \textbf{HiAP}：指在制定所有公共政策时系统性地考虑其对健康的影响，明确责任并加强跨部门合作，旨在改善人群健康和健康公平。
\end{enumerate}

\subsubsection*{三、简答题}
\begin{enumerate}[leftmargin=*, itemsep=3pt]
    \item \textbf{研究对象}：社会人群的健康。\textbf{研究内容}三大方面：①研究社会卫生状况（人群健康状况评价）；②研究影响人群健康的因素（特别是社会因素）；③研究社会卫生策略与措施。
    \item 五大任务：①倡导积极的健康观；②发现和评价社会卫生问题；③揭示社会因素与健康的因果关系；④制定社会卫生策略和措施；⑤培养社会医学专业人才。
    \item 特色理论：社会健康观、社会因素决定论、健康公平理论、整体健康观、大卫生观、健康社会决定因素理论。
    \item 重要阶段：①古代：中医"天人合一"整体观；②近代：陈志潜定县模式（1941年）；③新中国成立后：爱国卫生运动；④改革开放后：1984年首部教材出版；⑤21世纪：健康中国战略实施。
\end{enumerate}

\newpage
